\section{A tétel lényege és vizsgaszintű vázlata}

\begin{summarybox}
A 16. tétel az operációs rendszer folyamatkezelésének alapfogalmait tárgyalja. A folyamat egy program végrehajtás alatti példánya, amelyhez címtér, erőforrások, állapot és processzvezérlő blokk tartozik. A folyamat kontextusa teszi lehetővé, hogy az operációs rendszer megszakítsa, később pedig ugyanonnan folytassa a végrehajtást. A szál vagy fonál egy processzen belüli végrehajtási út, amely saját vezérlési állapottal rendelkezik, de a processz erőforrásain osztozik a többi szállal. A folyamat- és szálállapotok azt mutatják meg, hogy az adott végrehajtási egység CPU-t használ, CPU-ra vár, valamilyen eseményre vár, felfüggesztett állapotban van, vagy már befejeződött.
\end{summarybox}

Vizsgán célszerű a következő sorrendet követni:

\begin{enumerate}
  \item az operációs rendszer folyamatmenedzseri szerepe;
  \item program, processz, taszk és szál fogalmának elválasztása;
  \item a folyamat memóriaképe és a folyamat kontextusa;
  \item PCB, processztábla és a folyamat nyilvántartása;
  \item folyamatok létrehozása, életciklusa és befejezése;
  \item processzállapotok: egyszerű háromállapotú modell és bővített állapotmodell;
  \item állapotátmenetek: ütemezés, preempció, I/O-várakozás, esemény bekövetkezése, felfüggesztés;
  \item processz futási módok: felhasználói mód és kernel mód;
  \item szálak fogalma, előnyei, felhasználói és kernel szintű szálak;
  \item taszk- és fonálállapotok, tipikus állapotátmenetek.
\end{enumerate}

A legfontosabb vizsgamondat: a processz a program futó példánya saját címtérrel és PCB-ben tárolt kontextussal, a szál pedig a processzen belüli végrehajtási egység, amely saját programszámlálóval, regiszterállapottal és veremmel rendelkezik, miközben megosztja a processz kódját, adatait és erőforrásait.

\begin{center}
\begin{tabularx}{\textwidth}{L{30mm}X X}
\toprule
\textbf{Fogalom} & \textbf{Jelentése} & \textbf{Vizsgán fontos megjegyzés} \\
\midrule
Program & Háttértáron tárolt, statikus utasítássorozat és a hozzá tartozó adatok. & Önmagában még nem fut; akkor lesz belőle folyamat, ha az operációs rendszer betölti és végrehajtás alá helyezi. \\
Processz / folyamat & Egy program végrehajtás alatti példánya. & Saját címtérrel, vezérlési információval és erőforrásokkal rendelkezik; az operációs rendszer PCB-ben tartja nyilván. \\
Taszk & Feladat vagy végrehajtható egység; a pontos jelentése rendszerfüggő. & Sok környezetben közel áll a processzhez, más rendszerekben inkább az ütemezhető végrehajtási egységet jelenti. \\
Szál / fonál & Egy processzen belüli végrehajtási út. & Saját programszámlálóval, regiszterállapottal és veremmel rendelkezik, de megosztja a processz címtérét és erőforrásait. \\
Ütemezhető egység & Az az egység, amelynek az ütemező CPU-időt ad. & Modern rendszerekben gyakran a szál az ütemezhető egység, míg a processz inkább erőforrás-tulajdonos. \\
\bottomrule
\end{tabularx}

\end{center}

\section{Az operációs rendszer folyamatmenedzseri szerepe}

\subsection{Miért központi fogalom a folyamat?}

Az operációs rendszer egyik alapfeladata, hogy több program futását látszólag vagy ténylegesen párhuzamosan szervezze. Ehhez nyilván kell tartania, hogy melyik program milyen állapotban van, milyen erőforrásokat használ, mikor kaphat CPU-t, mire vár, és hogyan lehet biztonságosan megszakítani vagy folytatni.

A folyamatmenedzser feladatai közé tartozik:

\begin{itemize}
  \item folyamatok létrehozása és megszüntetése;
  \item CPU-idő kiosztása, vagyis ütemezés;
  \item folyamatállapotok nyilvántartása;
  \item kontextusváltás kezelése;
  \item folyamatok közötti kommunikáció és együttműködés támogatása;
  \item erőforrások, jogosultságok és védelem kezelése.
\end{itemize}

Egy egyprocesszoros rendszerben egy adott pillanatban csak egy végrehajtási egység fut ténylegesen a CPU-n, mégis több program tűnhet aktívnak. Ezt az operációs rendszer gyors váltásokkal, időszeletekkel, megszakításokkal és ütemezéssel éri el. Többmagos vagy többprocesszoros rendszerben már valódi párhuzamos végrehajtás is lehet.

\subsection{Program, folyamat és szoftver}

A \textbf{program} statikus fogalom: utasítások és adatok gyűjteménye, amely háttértáron található. Egy program akár hosszú ideig létezhet úgy, hogy éppen nem fut. A \textbf{szoftver} ennél tágabb fogalom: programok, könyvtárak, konfigurációk és egyéb adatok együttese.

A \textbf{folyamat} vagy \textbf{processz} dinamikus fogalom. Akkor jön létre, amikor a program végrehajtás alá kerül. Ugyanabból a programból egyszerre több folyamat is futhat, például két külön megnyitott böngészőablak vagy több terminálból indított azonos parancs. Ezek a folyamatok ugyanahhoz a programkódhoz kapcsolódhatnak, de eltérő címtérrel, állapottal, nyitott fájlokkal és végrehajtási pozícióval rendelkeznek.

Röviden:

\begin{itemize}
  \item a program \textbf{passzív}: háttértáron tárolt kód és adat;
  \item a processz \textbf{aktív}: futásban lévő programegyed;
  \item a processzhez állapot, erőforrások és végrehajtási kontextus tartozik.
\end{itemize}

\subsection{Taszk és szál viszonya}

A \textbf{taszk} szó jelentése nem minden rendszerben azonos. Egyes környezetekben szinte a processz szinonimája, máshol általánosabb értelemben feladatot vagy ütemezhető egységet jelent. Vizsgán érdemes jelezni, hogy a pontos jelentés rendszerfüggő. Általános operációs rendszeres szemléletben a taszk olyan végrehajtandó feladat, amelynek állapota és ütemezési helyzete van.

A \textbf{szál} vagy \textbf{fonál} a processzen belüli végrehajtási egység. Egy processz legalább egy szállal rendelkezik, de több szála is lehet. A szálak ugyanazon processz közös erőforrásait használják, ezért könnyebb közöttük adatot megosztani, ugyanakkor nagyobb figyelmet igényel a szinkronizáció.

\section{A processz felépítése és kontextusa}

\subsection{A folyamat memóriaképe}

Amikor egy program folyamatként elindul, az operációs rendszer memóriaterületet és végrehajtási környezetet rendel hozzá. A folyamat memóriaképe tipikusan több logikai részre bontható:

\begin{itemize}
  \item \textbf{text vagy kódterület}: a program végrehajtható utasításai;
  \item \textbf{data terület}: globális és statikus változók, inicializált és inicializálatlan adatok;
  \item \textbf{heap}: futás közben dinamikusan foglalt memória;
  \item \textbf{stack}: függvényhívások, lokális változók, visszatérési címek és paraméterek területe.
\end{itemize}

A processz nem csak memóriából áll. Hozzátartozhatnak nyitott fájlok, hálózati kapcsolatok, jogosultsági adatok, környezeti változók, ütemezési információk és különböző operációs rendszerbeli nyilvántartások is.

\subsection{Folyamatkontextus}

A \textbf{folyamatkontextus} azoknak az információknak az összessége, amelyek szükségesek ahhoz, hogy a folyamat végrehajtása megszakítható és később folytatható legyen. Ha az operációs rendszer elveszi a CPU-t egy folyamattól, el kell mentenie, hogy hol tartott a folyamat, milyen regiszterértékei voltak, milyen veremmutatót használt, milyen állapotban volt, és milyen erőforrások tartoztak hozzá.

A kontextus két nagy részre bontható:

\begin{itemize}
  \item \textbf{felhasználói szintű kontextus}: a folyamat saját címtérben látható adatai, kódja, vermének és heapjének tartalma;
  \item \textbf{kernel szintű kontextus}: az operációs rendszer által tárolt vezérlési adatok, például PCB, ütemezési információk, memóriakezelési adatok, I/O-állapotok.
\end{itemize}

A kontextus fogalma azért fontos, mert ettől válik lehetővé a többfeladatos működés. A folyamat szempontjából úgy tűnhet, mintha folyamatosan futna, miközben az operációs rendszer valójában többször megszakíthatja és később visszaadhatja neki a CPU-t.

\subsection{PCB és processztábla}

A \textbf{PCB} jelentése Process Control Block, vagyis processzvezérlő blokk. Ez az operációs rendszer egyik belső adatszerkezete, amely egy folyamat vezérléséhez és nyilvántartásához szükséges adatokat tartalmazza. A PCB nem a felhasználói program szabadon módosítható része, hanem a kernel által kezelt rendszeradat.

A processzek PCB-it az operációs rendszer rendszerint processztáblában vagy hasonló belső adatszerkezetekben tartja nyilván. Ezeknek mindig elérhetőnek kell lenniük a kernel számára, mert megszakítás, rendszerhívás, ütemezési döntés vagy I/O-esemény esetén gyorsan meg kell találni a megfelelő folyamat adatait.

\begin{center}
\begin{tabularx}{\textwidth}{L{36mm}X X}
\toprule
\textbf{Kontextuselem} & \textbf{Tartalma} & \textbf{Miért szükséges?} \\
\midrule
Azonosítók & Processzazonosító, szülőprocessz, tulajdonos, jogosultsági adatok. & Ezek alapján tudja az operációs rendszer megkülönböztetni és ellenőrizni a folyamatokat. \\
Állapot & Például új, futásra kész, futó, blokkolt, felfüggesztett vagy befejezett állapot. & Az ütemező ebből látja, hogy a folyamat kaphat-e CPU-t, vagy valamilyen eseményre vár. \\
Programsámláló és regiszterek & A következő utasítás címe, általános és speciális CPU-regiszterek, veremmutató. & Kontextusváltáskor ezeket menteni és visszaállítani kell, hogy a folyamat ott folytatódjon, ahol félbeszakadt. \\
Ütemezési adatok & Prioritás, CPU-használat, időszelet, várakozási idő, ütemezési sorokhoz tartozó hivatkozások. & Ezek alapján dönt az ütemező a futtatási sorrendről és a CPU-idő elosztásáról. \\
Memória-információ & Címtér, lap- vagy szegmenstábla adatai, memóriahatárok, jogosultságok. & A folyamat izolációjához, címleképzéséhez és védelméhez szükséges. \\
I/O és fájlinformáció & Nyitott fájlok, eszközök, I/O kérések, kommunikációs kapcsolatok. & A folyamat erőforrásait és függő műveleteit tartja nyilván. \\
Nyilvántartási adatok & Elhasznált CPU-idő, indulási idő, kvóták, statisztikák. & Teljesítményméréshez, elszámoláshoz és adminisztrációhoz használható. \\
\bottomrule
\end{tabularx}

\end{center}

\section{Folyamatok életciklusa}

\subsection{Folyamat létrehozása}

Új folyamat többféleképpen jöhet létre. Tipikus esetek:

\begin{itemize}
  \item rendszerindításkor az operációs rendszer alapfolyamatokat indít;
  \item egy futó folyamat új folyamatot hoz létre, például UNIX-szerű rendszerekben \code{fork()} segítségével;
  \item felhasználói kérésre elindul egy program;
  \item kötegelt vagy háttérfeladat végrehajtása kezdődik meg.
\end{itemize}

A létrehozás során az operációs rendszer azonosítót ad a folyamatnak, létrehozza a szükséges nyilvántartási adatokat, előkészíti a címtartományt, beállítja a kezdő végrehajtási pontot, és a folyamatot a megfelelő állapotba helyezi. A folyamat ezután rendszerint futásra kész állapotba kerül, vagyis vár arra, hogy az ütemező CPU-t rendeljen hozzá.

UNIX-szerű rendszerekben gyakori a szülő--gyermek folyamatkapcsolat. A gyermekfolyamatot egy már létező folyamat hozza létre, ezért folyamatfa vagy folyamathierarchia alakulhat ki. Háttérben futó, szolgáltatásszerű folyamatokat gyakran démonfolyamatoknak nevezünk.

\subsection{Folyamat befejezése}

A folyamat több okból fejeződhet be:

\begin{itemize}
  \item \textbf{normál kilépés}: a folyamat elvégezte a feladatát;
  \item \textbf{kezelt hiba}: olyan hiba történt, amely után a folyamat már nem tud értelmesen tovább futni;
  \item \textbf{végzetes hiba}: például hibás memóriahivatkozás vagy nem kezelhető kivételes helyzet;
  \item \textbf{más folyamat vagy felhasználó beavatkozása}: például \code{kill()} jellegű művelet.
\end{itemize}

Befejezéskor az operációs rendszernek takarítási feladata van. Be kell állítani a visszatérési állapotot, kezelni kell a gyermekfolyamatokat, fel kell szabadítani a memóriát és az I/O-erőforrásokat, majd törölni kell a folyamat nyilvántartási adatait, ha már nincs rájuk szükség.

\subsection{Folyamatok együttélése}

A folyamatok a rendszer erőforrásaiért versenyeznek. CPU-időt, memóriát, háttértárat, I/O-eszközöket, fájlokat vagy hálózati kapcsolatokat használhatnak. Emellett kommunikálhatnak is egymással, például csővezetéken, üzenetsoron, osztott memórián vagy egyéb IPC-mechanizmuson keresztül.

A versengés és együttműködés miatt van szükség ütemezésre, szinkronizációra, jogosultságkezelésre és erőforrás-védelemre. A folyamatállapotok azt mutatják, hogy egy adott folyamat éppen hol tart ebben a versenyben.

\section{Processzállapotok és állapotátmenetek}

\subsection{CPU-szakasz és I/O-szakasz}

A folyamat végrehajtása tipikusan CPU-igényes és I/O-igényes szakaszok váltakozásából áll. Egy ideig számítás történik, majd a folyamat például fájlra, billentyűzetre, hálózatra vagy más eszközre vár. Amíg a folyamat I/O-ra vár, nem érdemes CPU-t adni neki, mert nem tudna hasznos munkát végezni.

Ez indokolja a multiprogramozást. Amikor az egyik folyamat blokkolódik, az operációs rendszer másik futásra kész folyamatot választhat. Így a CPU kihasználtsága javul, miközben a folyamatok felváltva haladnak.

\subsection{Egyszerű háromállapotú modell}

A legegyszerűbb működő modell három állapotot különböztet meg:

\begin{itemize}
  \item \textbf{Running / futó}: a folyamat éppen CPU-t használ;
  \item \textbf{Ready / futásra kész}: a folyamat futtatható, csak CPU-ra vár;
  \item \textbf{Blocked / blokkolt}: a folyamat valamilyen eseményre vár, ezért nem futtatható.
\end{itemize}

A legfontosabb átmenetek:

\begin{itemize}
  \item futó folyamat I/O-kérést ad ki: \textbf{running $\rightarrow$ blocked};
  \item az ütemező elveszi tőle a CPU-t: \textbf{running $\rightarrow$ ready};
  \item az ütemező kiválasztja: \textbf{ready $\rightarrow$ running};
  \item az I/O vagy várt esemény befejeződik: \textbf{blocked $\rightarrow$ ready}.
\end{itemize}

\subsection{Bővített állapotmodell}

A valós operációs rendszerekben a háromállapotú modell gyakran kiegészül további állapotokkal. Ilyen például az új állapot, a befejezett állapot, illetve a felfüggesztett állapotok. Felfüggesztésre akkor lehet szükség, ha a rendszert tehermentesíteni kell, ha memóriát kell felszabadítani, vagy ha adminisztratív okból átmenetileg meg kell állítani egy folyamatot.

\begin{center}
\small
\begin{tabularx}{\textwidth}{L{32mm}X X}
\toprule
\textbf{Állapot} & \textbf{Jelentése} & \textbf{Tipikus átmenet} \\
\midrule
Új & A folyamat létrejött, de még nem feltétlenül került be a futásra kész sorba. & Létrehozás után felvétel a ready sorba. \\
Futásra kész & A folyamat futtatható, csak CPU-ra vár. & Az ütemező kiválasztja: ready $\rightarrow$ running. \\
Futó & A folyamat éppen CPU-t használ. & Időszelet lejár: running $\rightarrow$ ready; I/O kérés: running $\rightarrow$ blocked; befejezés: running $\rightarrow$ terminated. \\
Blokkolt / várakozó & A folyamat nem tud futni, mert valamilyen eseményre vár, például I/O befejezésére. & Az esemény bekövetkezik: blocked $\rightarrow$ ready. \\
Felfüggesztett ready & A folyamat futtatható lenne, de átmenetileg ki van vonva a CPU-ért folyó versenyből, gyakran memóriakezelési okból. & Resume után visszakerül a ready állapotba. \\
Felfüggesztett blokkolt & A folyamat eseményre vár és közben felfüggesztett állapotban van. & Az esemény bekövetkezhet, de a folyamat csak resume után lehet futásra kész. \\
Befejezett & A folyamat végrehajtása lezárult vagy megszakadt. & Az operációs rendszer felszabadítja az erőforrásokat és törli a nyilvántartási adatokat. \\
\bottomrule
\end{tabularx}

\end{center}

\begin{figure}[htbp]
\centering
\resizebox{0.92\textwidth}{!}{%
\begin{tikzpicture}[font=\small]
  \node[zvflowprocess, fill=black!5, minimum width=20mm] (new) at (0,0) {Új};
  \node[zvflowprocess, fill=lightaccent, minimum width=31mm] (ready) at (3.1,0) {Futásra kész};
  \node[zvflowprocess, fill=warnbg, minimum width=23mm] (running) at (6.4,0) {Futó};
  \node[zvflowprocess, fill=black!5, minimum width=27mm] (terminated) at (9.5,0) {Befejezett};

  \node[zvflowprocess, fill=lightaccent!65, minimum width=29mm] (blocked) at (6.4,-2.0) {Blokkolt};
  \node[zvflowprocess, fill=lightaccent!45, minimum width=30mm, align=center] (sready) at (1.9,-4.0) {Felfüggesztett\\ ready};
  \node[zvflowprocess, fill=lightaccent!45, minimum width=31mm, align=center] (sblocked) at (7.0,-4.0) {Felfüggesztett\\ blokkolt};

  \draw[arrow] (new) -- (ready);
  \draw[arrow, bend left=25] (ready) to (running);
  \draw[arrow, bend left=25] (running) to (ready);
  \draw[arrow] (running) -- (terminated);
  \draw[arrow] (running) -- (blocked);
  \draw[arrow] (blocked.west) to[out=180,in=-60] (ready.south);
  \draw[arrow] (ready.south west) to[out=-110,in=80] (sready.north);
  \draw[arrow] (sready.north east) to[out=80,in=-120] (ready.south);
  \draw[arrow] (blocked) -- (sblocked);
  \draw[arrow] (sblocked) -- (blocked);
  \draw[arrow] (sblocked) -- (sready);
\end{tikzpicture}%
}
\caption{Folyamatállapotok bővített modellje és tipikus állapotátmenetei}
\label{fig:zv16_processz_allapotmodell}
\end{figure}


A bővített modellben fontos megkülönböztetni a \textbf{blokkolt} és a \textbf{felfüggesztett} állapotot. A blokkolt folyamat azért nem fut, mert valamilyen eseményre vár. A felfüggesztett folyamatot viszont az operációs rendszer vagy a felhasználó vonta ki átmenetileg az aktív versenyből. Egy folyamat lehet egyszerre eseményre váró és felfüggesztett is.

\section{Kontextusváltás és ütemezés}

\subsection{Mikor van szükség ütemezésre?}

Ütemezésre akkor van szükség, amikor dönteni kell, hogy melyik futásra kész folyamat vagy szál kapja meg a CPU-t. Ilyen döntési helyzet például:

\begin{itemize}
  \item a futó folyamat I/O-ra vár és blokkolódik;
  \item lejár a futó folyamat időszelete;
  \item egy blokkolt folyamat eseménye befejeződik, ezért futásra kész lesz;
  \item egy folyamat befejeződik;
  \item új, magasabb prioritású feladat válik futtathatóvá.
\end{itemize}

Ha az operációs rendszer csak akkor vált, amikor a futó folyamat önként lemond a CPU-ról vagy blokkolódik, akkor \textbf{nem preemptív} ütemezésről beszélünk. Ha az operációs rendszer például órajel-megszakítás alapján kényszerítetten is elveheti a CPU-t, akkor \textbf{preemptív} ütemezésről van szó.

\subsection{Kontextusváltás menete}

Kontextusváltáskor az operációs rendszer elmenti a jelenleg futó folyamat vagy szál végrehajtási állapotát, majd visszaállítja a következő kiválasztott egység állapotát. A mentés és visszaállítás tartalmazhatja a programszámlálót, regisztereket, veremmutatót, processzállapotot, memóriakezelési információkat és ütemezési adatokat.

\begin{figure}[htbp]
\centering
\resizebox{0.95\textwidth}{!}{%
\begin{tikzpicture}[font=\scriptsize]
  \node[zvflowprocess, fill=warnbg, minimum width=18mm, minimum height=8mm] (p1run) at (0,0) {P1 fut};
  \node[zvflowprocess, fill=black!6, minimum width=25mm, minimum height=8mm] (save) at (2.7,0) {P1 mentése};
  \node[zvflowprocess, fill=lightaccent, minimum width=22mm, minimum height=8mm] (sched) at (5.3,0) {ütemező};
  \node[zvflowprocess, fill=black!6, minimum width=27mm, minimum height=8mm] (restore) at (8.0,0) {P2 visszaállítás};
  \node[zvflowprocess, fill=warnbg, minimum width=18mm, minimum height=8mm] (p2run) at (10.8,0) {P2 fut};

  \draw[arrow] (p1run) -- node[above, font=\tiny] {megszakítás} (save);
  \draw[arrow] (save) -- (sched);
  \draw[arrow] (sched) -- (restore);
  \draw[arrow] (restore) -- node[above, font=\tiny] {visszatérés} (p2run);
\end{tikzpicture}%
}
\caption{Kontextusváltás leegyszerűsített menete két folyamat között}
\label{fig:zv16_kontextusvaltas}
\end{figure}


A kontextusváltás szükséges a többfeladatos működéshez, de önmagában nem végez felhasználói szempontból hasznos munkát. Ezért az operációs rendszerek törekednek arra, hogy a váltás elég gyors legyen, és ne történjen indokolatlanul gyakran.

\section{Processz futási módok}

\subsection{Felhasználói mód és kernel mód}

A modern processzorok és operációs rendszerek rendszerint megkülönböztetnek \textbf{felhasználói módot} és \textbf{kernel módot}. Erre azért van szükség, mert a felhasználói programok nem férhetnek hozzá közvetlenül minden hardveres és rendszerbeli erőforráshoz. Ha ezt megtehetnék, egy hibás vagy rosszindulatú program könnyen tönkretehetné más folyamatok adatait vagy az egész rendszert.

\begin{itemize}
  \item \textbf{Felhasználói mód}: a folyamat korlátozott jogosultságokkal fut. Nem hajthat végre privilegizált utasításokat, és közvetlenül nem kezelheti a hardver nagy részét.
  \item \textbf{Kernel mód}: az operációs rendszer magja teljesebb jogosultságokkal fut. Itt hajthatók végre a memóriakezeléshez, I/O-hoz, megszakításkezeléshez és ütemezéshez szükséges műveletek.
\end{itemize}

A folyamat normál programkódja felhasználói módban fut. Amikor rendszerhívást végez, például fájlt olvas, memóriát kér vagy folyamatot hoz létre, vezérelt módváltás történik kernel módba. Megszakítás vagy kivétel esetén szintén a kernel veszi át a vezérlést.

\subsection{Módváltás és védelem}

A módváltás nem ugyanaz, mint a processzváltás. Rendszerhíváskor gyakran ugyanannak a folyamatnak a nevében fut kernelkód, csak magasabb jogosultsági szinten. Kontextusváltáskor viszont az operációs rendszer másik folyamat vagy szál végrehajtási állapotára vált át.

A futási módok lényege a védelem:

\begin{itemize}
  \item a felhasználói program ne módosíthassa közvetlenül a kernel memóriáját;
  \item a folyamatok címtartományai legyenek elválasztva;
  \item az I/O-eszközök kezelése ellenőrzött rendszerhívásokon keresztül történjen;
  \item a hibás program ne állíthassa le közvetlenül az egész rendszert.
\end{itemize}

\section{Szálak és fonalak}

\subsection{A szál fogalma}

A \textbf{szál} vagy \textbf{fonál} egy processzen belül létrehozott végrehajtási egység. A processz erőforrás-tulajdonosként is felfogható, a szál pedig végrehajtási útvonalnak tekinthető. Egy processz szálai ugyanazt a címtartományt használják, de mindegyiknek saját végrehajtási állapota van.

Egy szálhoz tipikusan tartozik:

\begin{itemize}
  \item saját programszámláló;
  \item saját CPU-regiszterállapot;
  \item saját verem;
  \item saját szálazonosító és ütemezési állapot.
\end{itemize}

A szálak viszont megosztják:

\begin{itemize}
  \item a processz kódterületét;
  \item a globális adatokat és a heapet;
  \item a nyitott fájlokat és egyéb processzszintű erőforrásokat;
  \item a processz jogosultsági és címtérbeli kereteit.
\end{itemize}

\begin{figure}[htbp]
\centering
\begin{tikzpicture}[font=\small]
  \node[draw, rounded corners, fill=lightaccent, minimum width=12.8cm, minimum height=5.8cm] (proc) at (0,0) {};
  \node[anchor=west, font=\bfseries\sffamily\large, color=accent] at (-5.85,2.25) {Processz};

  \node[zvflowprocess, fill=black!5, minimum width=29mm, minimum height=8mm] (code) at (-4.35,1.2) {kód / text};
  \node[zvflowprocess, fill=black!5, minimum width=25mm, minimum height=8mm] (data) at (-1.6,1.2) {adatok};
  \node[zvflowprocess, fill=black!5, minimum width=25mm, minimum height=8mm] (heap) at (1.1,1.2) {heap};
  \node[zvflowprocess, fill=black!5, minimum width=31mm, minimum height=8mm] (files) at (4.2,1.2) {nyitott fájlok};

  \draw[decorate, decoration={brace, mirror, amplitude=5pt}, thick]
    (-5.2,0.58) -- (5.2,0.58)
    node[midway, below=7pt, align=center] {megosztott processzszintű erőforrások és közös címtér};

  \node[draw, rounded corners, fill=warnbg, minimum width=34mm, minimum height=17mm, align=center] (t1) at (-3.6,-1.4) {Szál 1\\{\scriptsize PC + regiszterek}\\{\scriptsize verem}};
  \node[draw, rounded corners, fill=warnbg, minimum width=34mm, minimum height=17mm, align=center] (t2) at (0,-1.4) {Szál 2\\{\scriptsize PC + regiszterek}\\{\scriptsize verem}};
  \node[draw, rounded corners, fill=warnbg, minimum width=34mm, minimum height=17mm, align=center] (t3) at (3.6,-1.4) {Szál 3\\{\scriptsize PC + regiszterek}\\{\scriptsize verem}};

  \draw[arrow] (-3.6,0.55) -- (t1.north);
  \draw[arrow] (0,0.55) -- (t2.north);
  \draw[arrow] (3.6,0.55) -- (t3.north);

  \node[font=\scriptsize, align=center] at (0,-2.55) {Minden szálnak saját végrehajtási állapota van, miközben a processz erőforrásait közösen használják.};
\end{tikzpicture}
\caption{Processz és szál kapcsolata: közös erőforrások, saját végrehajtási állapot}
\label{fig:zv16_processz_es_szal_eroforrasok}
\end{figure}


\subsection{Miért előnyösek a szálak?}

A szálak használatának több gyakorlati előnye van:

\begin{itemize}
  \item \textbf{gyorsabb reakció}: ha egy szál I/O-ra vár, a processz más szálai még dolgozhatnak;
  \item \textbf{olcsóbb vezérlés}: szálat létrehozni és szálak között váltani általában olcsóbb, mint teljes processzek között;
  \item \textbf{erőforrás-megosztás}: a szálak közös memóriát és fájlokat használhatnak;
  \item \textbf{párhuzamosság}: többmagos rendszeren több szál ténylegesen párhuzamosan futhat;
  \item \textbf{természetes feladatfelbontás}: például egy médialejátszóban külön szál kezelheti a hangot, a videót, a hálózatot és a felhasználói felületet.
\end{itemize}

A közös címtér miatt ugyanakkor a szálak veszélyesebbek is lehetnek: ha több szál ugyanazt az adatot módosítja, szinkronizáció nélkül versenyhelyzet, inkonzisztens állapot vagy holtpont alakulhat ki.

\begin{center}
\begin{tabularx}{\textwidth}{L{32mm}X X}
\toprule
\textbf{Szempont} & \textbf{Processz} & \textbf{Szál / fonál} \\
\midrule
Alapjelentés & Program végrehajtás alatti példánya. & Végrehajtási út egy processzen belül. \\
Címtér & Általában saját, védett címtérrel rendelkezik. & Megosztja a processz címtérét a többi szállal. \\
Saját adatok & PCB, címtér, nyitott erőforrások, legalább egy végrehajtási szál. & Saját programszámláló, regiszterállapot és verem. \\
Erőforrás-megosztás & Processzek között drágább és szabályozottabb, például IPC kellhet. & Azonos processzen belüli szálak könnyen megosztják a kódot, adatokat, heapet és fájlokat. \\
Létrehozás és váltás költsége & Általában drágább, mert több erőforrás és több védelmi információ kapcsolódik hozzá. & Általában olcsóbb, ezért nevezik könnyűsúlyú folyamatnak is. \\
Hibatűrés & Egy processz hibája gyakran elszigetelhető más processzektől. & Egy szál hibája az egész processzt érintheti, mert közös címtéren osztoznak. \\
\bottomrule
\end{tabularx}

\end{center}

\subsection{Felhasználói és kernel szintű szálak}

A szálak megvalósítása alapján két alapvető eset különíthető el. \textbf{Felhasználói szintű szálaknál} a szálkezelést egy felhasználói könyvtár végzi, ezért az operációs rendszer gyakran csak a teljes processzt látja. \textbf{Kernel szintű szálaknál} az operációs rendszer ismeri és külön ütemezi a szálakat.

\begin{center}
\begin{tabularx}{\textwidth}{L{34mm}X X}
\toprule
\textbf{Típus} & \textbf{Jellemző} & \textbf{Következmény} \\
\midrule
Felhasználói szintű szál & A szálakat felhasználói könyvtár kezeli, az operációs rendszer csak a processzt látja. & A szálváltás gyors lehet, de ha egy szál blokkoló rendszerhívást végez, az egész processz blokkolódhat. \\
Kernel szintű szál & A szálakat az operációs rendszer kernelje ismeri és ütemezi. & Egy szál blokkolása nem feltétlenül állítja meg a processz többi szálát; a váltás viszont drágább lehet. \\
Egy--több modell & Több felhasználói szál tartozik egy kernel végrehajtási egységhez. & Egyszerű, de valódi többmagos párhuzamosságot korlátozottan ad. \\
Egy--egy modell & Minden felhasználói szálhoz külön kernel szál tartozik. & Jó párhuzamosságot ad, de sok kernel szál erőforrásigényes lehet. \\
Több--több modell & Több felhasználói szál több kernel szálra képeződik le. & Rugalmasabb, de bonyolultabb megvalósítani. \\
\bottomrule
\end{tabularx}

\end{center}

Ha a szálak csak felhasználói szinten léteznek, akkor egy blokkoló rendszerhívás az egész processzt megállíthatja, mert az operációs rendszer nem látja a processzen belüli többi futtatható szálat. Kernel szintű szálaknál az operációs rendszer egy másik szálat is futtathat ugyanabból a processzből, amíg az egyik szál várakozik.

\section{Taszk- és fonálállapotok}

\subsection{Taszkállapotok}

A taszk állapotmodellje általában hasonló a processz állapotmodelljéhez, mert a taszk is végrehajtandó vagy ütemezendő egységet jelent. A leggyakoribb állapotok:

\begin{itemize}
  \item \textbf{created / new}: a taszk létrejött;
  \item \textbf{ready}: futtatható, CPU-ra vár;
  \item \textbf{running}: éppen végrehajtás alatt van;
  \item \textbf{waiting / blocked}: eseményre, erőforrásra vagy I/O-ra vár;
  \item \textbf{suspended}: átmenetileg kivonták az ütemezésből;
  \item \textbf{terminated}: befejeződött.
\end{itemize}

Az átmenetek oka lehet indítás, ütemezői kiválasztás, időszelet lejárta, várakozás kezdete, esemény bekövetkezése, felfüggesztés, folytatás vagy befejezés. Valós idejű operációs rendszerekben a taszk fogalom különösen gyakori, mert ott a feladatok határidővel, prioritással és eseményekre adott reakcióval kapcsolódnak össze.

\subsection{Fonálállapotok}

A fonál vagy szál állapotai a processzállapotokhoz hasonlítanak, de az állapot a processzen belüli végrehajtási útra vonatkozik. Egy többszálú processzben elképzelhető, hogy az egyik szál blokkolt, miközben egy másik szál ugyanazon processzen belül futásra kész vagy futó állapotban van.

Tipikus fonálállapotok:

\begin{itemize}
  \item \textbf{új}: a szál objektuma vagy leírója létrejött, de még nem indult el;
  \item \textbf{futásra kész}: a szál végrehajtható, CPU-ra vár;
  \item \textbf{futó}: a szál CPU-n van;
  \item \textbf{blokkolt vagy várakozó}: zárolásra, I/O-ra, időzítőre vagy másik szál jelzésére vár;
  \item \textbf{befejezett}: a szál futása véget ért.
\end{itemize}

A fontos átmenetek:

\begin{itemize}
  \item indítás után a szál futásra kész lesz;
  \item az ütemező futásra kész szálat választ, így futó állapotba kerül;
  \item időszelet lejártakor vagy magasabb prioritású szál megjelenésekor visszakerülhet futásra kész állapotba;
  \item I/O, zárolás vagy várakozás miatt blokkolódhat;
  \item a várt esemény bekövetkezése után ismét futásra kész lesz;
  \item a szálfüggvény végén befejezett állapotba kerül.
\end{itemize}

\begin{warningbox}
Vizsgán érdemes külön kiemelni: ha a szálak kernel szintűek, akkor egy szál blokkolódása nem feltétlenül blokkolja az egész processzt. Ha viszont a szálakat csak felhasználói szintű könyvtár kezeli, akkor egy blokkoló rendszerhívás az egész processzt megakaszthatja.
\end{warningbox}

\section{Gyakori vizsgahibák és összefoglalás}

\subsection{Gyakori félreértések}

\begin{itemize}
  \item \textbf{A program és a processz nem ugyanaz.} A program statikus, a processz futó példány.
  \item \textbf{A blokkolt folyamat nem azért nem fut, mert nincs kiválasztva.} A blokkolt folyamat valamilyen eseményre vár, ezért akkor sem lenne futtatható, ha lenne szabad CPU.
  \item \textbf{A ready állapot nem jelent futást.} A folyamat készen áll a futásra, de még vár a CPU-ra.
  \item \textbf{A módváltás nem azonos a kontextusváltással.} Kernel módba át lehet lépni ugyanannak a folyamatnak a nevében is.
  \item \textbf{A szál nem teljesen önálló processz.} Saját végrehajtási állapota van, de sok erőforrást megoszt a processz többi szálával.
  \item \textbf{A szálak közötti közös memória előny és veszély is.} Gyors kommunikációt ad, de szinkronizáció nélkül versenyhelyzetet okozhat.
\end{itemize}

\subsection{Rövid vizsgaválasz-minta}

\begin{summarybox}
A processz egy program végrehajtás alatti példánya, amelyet az operációs rendszer hoz létre, ütemez, nyilvántart és szüntet meg. A processzhez saját címtér, erőforrások és PCB tartozik. A PCB tartalmazza többek között a processz azonosítóját, állapotát, programszámlálóját, regisztereit, ütemezési, memória- és I/O-adatait. A folyamat kontextusa teszi lehetővé a kontextusváltást: az operációs rendszer elmenti a futó folyamat állapotát, majd egy másik folyamat vagy szál állapotát állítja vissza. A fő folyamatállapotok: új, futásra kész, futó, blokkolt, felfüggesztett és befejezett. A szál a processzen belüli végrehajtási egység, amely saját programszámlálóval, regiszterekkel és veremmel rendelkezik, de megosztja a processz kódját, adatait, heapjét és nyitott erőforrásait. A futási módok közül felhasználói módban korlátozott jogosultságú programkód fut, kernel módban pedig az operációs rendszer privilegizált műveletei hajtódnak végre.
\end{summarybox}
